\documentclass[11pt,a4paper]{ctexart}
\usepackage[utf8]{inputenc}
\usepackage{geometry}
\geometry{left=2.5cm,right=2.5cm,top=2.8cm,bottom=2.8cm}
\usepackage{amsmath,amssymb,amsfonts}
\usepackage{enumitem}
\usepackage{titlesec}
\usepackage{xcolor}
\usepackage{tcolorbox}
\usepackage{hyperref}

% 引入绝对位置放置宏包用于生成水印及插入图片宏包
\usepackage{graphicx}
\usepackage{eso-pic}

% 解决带圈数字在部分编译器下报错的问题


% --- 颜色定义 ---
\definecolor{mainblue}{RGB}{26, 82, 118}
\definecolor{secblue}{RGB}{41, 128, 185}
\definecolor{bglight}{RGB}{244, 246, 247}
\definecolor{textdark}{RGB}{44, 62, 80}

% --- 页面美化与字体设置 ---
\hypersetup{colorlinks=true, linkcolor=mainblue, urlcolor=secblue}
\setlist[itemize]{leftmargin=2em, itemsep=0.3em}
\setlist[enumerate]{leftmargin=2em, itemsep=0.3em}

% --- 标题样式自定义 ---
\titleformat{\section}{\large\bfseries\color{mainblue}}{\thesection}{1em}{}[{\color{mainblue}\titlerule[1pt]}]
\titleformat{\subsection}{\normalsize\bfseries\color{secblue}}{\thesubsection}{1em}{}
\titleformat{\subsubsection}{\small\bfseries\color{textdark}}{\thesubsubsection}{1em}{}

% --- 水印设置 (每页三个，斜45度，居中分布) ---
\newcommand{\WMText}{贺昱琪学习资料}
\newcommand{\WMScale}{2.28}
\newcommand{\WMAngle}{45}
\colorlet{WMColor}{black!8}

\AddToShipoutPictureBG{%
  \AtPageLowerLeft{%
    \put(\LenToUnit{0.66\paperwidth},\LenToUnit{0.70\paperheight}){%
      \rotatebox{\WMAngle}{\scalebox{\WMScale}{\textcolor{WMColor}{\WMText}}}%
    }%
    \put(\LenToUnit{0.33\paperwidth},\LenToUnit{0.50\paperheight}){%
      \rotatebox{\WMAngle}{\scalebox{\WMScale}{\textcolor{WMColor}{\WMText}}}%
    }%
    \put(\LenToUnit{0.66\paperwidth},\LenToUnit{0.30\paperheight}){%
      \rotatebox{\WMAngle}{\scalebox{\WMScale}{\textcolor{WMColor}{\WMText}}}%
    }%
  }%
}

\begin{document}

\begin{center}
    \LARGE \textbf{贺昱琪专项训练三}
\end{center}

\vspace{0.5cm}

\noindent \textbf{一、认真填一填（每小题 3 分，共 30 分）}

\begin{enumerate}
    \item 若 $\frac{11}{m}$ 是最简真分数，已知 $m$ 是两位自然数，且 $m$ 的个位数字比十位数字大 5，且 $\frac{11}{m}$ 是纯循环小数，则满足以上条件的 $m=$\underline{\hspace{2cm}}。\vspace{1.2cm}
    
    \item 有四个自然数 $1, a, b, c$，满足 $1 < a < b < c$，这四个自然数两两求和可得到 6 个不同的数。把这 6 个数从小到大排列，求相邻两数的差，又得到 5 个数。这 5 个数恰好相等，则 $c = $\underline{\hspace{2cm}}。\vspace{1.2cm}
    
    \item 已知一个直角三角形的三边长为 6cm，8cm，10cm，则该三角形的高为\underline{\hspace{2cm}}。\vspace{1.2cm}
    
    \item 
    \begin{minipage}{0.65\textwidth}
        如图，小正方形的 $\frac{3}{4}$ 被阴影部分覆盖，大正方形的 $\frac{6}{7}$ 被阴影部分覆盖。小正方形的阴影部分面积与大正方形的阴影部分面积之比是\underline{\hspace{2cm}}。
    \end{minipage}
    \hfill
    \begin{minipage}{0.3\textwidth}
        \centering
        \includegraphics[width=\linewidth]{4.png}
    \end{minipage}
    \vspace{1.2cm}
    
    \item 
    \begin{minipage}{0.65\textwidth}
        如图，在 $\triangle ABC$ 中，$A_1, A_2, A_3, \dots, A_n$ 为 $BC$ 边上不同的 $n$ 个点。首先连接 $AA_1$，图中出现 3 个不同的三角形，再连接 $AA_2$，图中有 6 个不同的三角形；再连接 $AA_3$，图中有 10 个不同的三角形，$\dots$。若一直连接到 $A_n$，则图中共有不同的三角形\underline{\hspace{2cm}}个。
    \end{minipage}
    \hfill
    \begin{minipage}{0.3\textwidth}
        \centering
        \includegraphics[width=\linewidth]{5.png}
    \end{minipage}
    \vspace{1.2cm}
    
    \item 
    \begin{minipage}{0.65\textwidth}
        我们通常把四条边都相等的四边形叫做菱形。探究菱形的两条对角线互相垂直且互相平分。如图，已知菱形 $ABCD$ 的边长为 5cm，对角线 $AC=6cm$，$BD=8cm$。对角线 $BD$ 上有一个动点 $P$，过点 $P$ 分别向边 $AB, AD$ 作垂线段 $PM, PN$，则 $PM+PN$ 有\underline{\hspace{2cm}}（填“最大、最小、定”）值，这个值为\underline{\hspace{2cm}}。
    \end{minipage}
    \hfill
    \begin{minipage}{0.3\textwidth}
        \centering
        \includegraphics[width=\linewidth]{6.png}
    \end{minipage}
    \vspace{1.2cm}
    
    \item 如果 $m-(m+1)+(m+2)$ 没有进位，我们叫 $n$ 为可连数。如 $32+33+34$ 没有进位，32 就是可连数，则 200 以内有\underline{\hspace{2cm}}个可连数。\vspace{1.2cm}
    
    \item 
    \begin{minipage}{0.65\textwidth}
        利用两块长方体木块测量一张桌子的高度。首先按图①方式放置，再交换两木块的位置，按图②方式放置。测量的数据如图，则桌子的高度是\underline{\hspace{2cm}}cm。
    \end{minipage}
    \hfill
    \begin{minipage}{0.3\textwidth}
        \centering
        \includegraphics[width=\linewidth]{8.png}
    \end{minipage}
    \vspace{1.2cm}
    
    \item 
    \begin{minipage}{0.65\textwidth}
        有一个深 4 分米的长方体容器，其内侧底面为边长 3 分米的正方形，当容器如右图所示倾斜时，容器内的水刚好不溢出。容器内的水有\underline{\hspace{2cm}}升。
    \end{minipage}
    \hfill
    \begin{minipage}{0.3\textwidth}
        \centering
        \includegraphics[width=\linewidth]{9.png}
    \end{minipage}
    \vspace{1.2cm}
    
    \item 若一个正整数能表示为两个正整数的平方差，则称这个正整数为“智慧数”（如 $3=2^2-1^2, 5=3^2-2^2$），已知智慧数按从小到大顺序构成如下数列：$3, 5, 7, 8, 9, 11, 12, 13, 15, 16, 17, 19, 20, 21, 23, 24, 25, \dots$，则第 100 个智慧数是\underline{\hspace{2cm}}。\vspace{1.2cm}
\end{enumerate}

\vspace{0.8cm}

\noindent \textbf{二、细心算一算}

\begin{enumerate}
    \item[11.] 计算（每小题 4 分，共 16 分）
    \begin{enumerate}
        \item[(1)] $\left[ 2\frac{4}{5} + \left(3\frac{1}{2} - 2.75\right) \times 1\frac{3}{5} \right] \div 3\frac{1}{5}$
        \vspace{3.5cm}
        
        \item[(2)] $\frac{4}{9} \times 0.8 + 2\frac{3}{4} \times 15 + 0.8 \times \frac{5}{9} - 1\frac{3}{4} \times 15$
        \vspace{3.5cm}
        
        \item[(3)] $138 \times 1\frac{1}{4} + 23.3 \times (2 - 75\%) - 125\% \times 38 + (1 + 0.25) \times 28.7$
        \vspace{3.5cm}
        
        \item[(4)] $\left( 2 + \frac{1}{2} + \frac{1}{3} + \dots + \frac{1}{9} \right) \times \left( \frac{1}{2} + \frac{1}{3} + \dots + \frac{1}{10} \right) - \left( 2 + \frac{1}{2} + \frac{1}{3} + \dots + \frac{1}{10} \right) \times \left( \frac{1}{2} + \frac{1}{3} + \dots + \frac{1}{9} \right)$
        \vspace{3.5cm}
    \end{enumerate}
\end{enumerate}

\vspace{0.8cm}

\noindent \textbf{三、用心想一想}

\begin{enumerate}
    \item[12.] (7分) 六年级甲、乙、丙三个班植树，任务分配是：甲班要植三个班总棵数的 $40\%$，乙、丙两班植树棵数的比是 $4:3$，当甲班植了 200 棵树时，正好完成三个班植树总棵数的 $\frac{2}{7}$。求丙班植树多少棵？
    \vspace{3.5cm}

    \item[13.] (8分) 某单位需要按中通或申通快递方式寄五所学校各寄一封信，这五封信的重量分别是 72g、90g、215g、340g、400g。根据两家快递公司提供的收费范围，在邮局查得相关资费标准如下：
    
    \begin{table}[htbp]
        \centering
        \begin{tabular}{|c|c|c|c|c|}
            \hline
            业务种类 & 计费单位 & 资费标准（元/件） & 挂号费（元/件） & 特制信封（元/个） \\
            \hline
            微小信 & 首重 100g 内，每重 20g & 0.8 & 3 & 0.5 \\
            \hline
            微小信 & 续重 101~2000g，每重 100g & 2.00 & 3 & 0.5 \\
            \hline
            特快专递 & 首重 1000g 内 & 5.00 & 3 & 1.0 \\
            \hline
        \end{tabular}
    \end{table}
    
    \begin{enumerate}
        \item[(1)] 重量为 90g 的信以“微小信”方式寄出，邮寄费为多少元？若以“特快专递”方式寄出呢？
        \vspace{2.5cm}
        \item[(2)] 这五封信分别以怎样的方式寄出最合算？请说明理由。
        \vspace{2.5cm}
        \item[(3)] 通过解答上述问题，你有什么启示？（请你用一、两句话说明）
        \vspace{2.5cm}
    \end{enumerate}

    \item[14.] (8分) 甲乙丙三人的行走速度分别为每分钟 80 米、60 米、50 米。甲、乙两人从 A 地，丙一人从 B 地同时相向出发，已知 AB 两地的距离是 300 米。如果甲乙从 A 往 B，丙往 A 的方向与甲乙相向，那么多少分钟后甲与丙，乙与丙的距离第一次相等？
    \vspace{3.5cm}
    
    \item[15.] (9分) 如图1是甲、乙两个圆柱形水槽的轴截面示意图，乙槽中有一圆柱形铁块放其中（圆柱形铁块的下底面完全落在水槽底面上）。现将甲槽中的水匀速注入乙槽。甲、乙两个水槽中水的深度 $y$（厘米）与注水时间 $x$（分钟）之间的关系如图2所示。根据图象提供的信息，解答下列问题：
    
    \begin{center}
        \includegraphics[width=0.6\textwidth]{15.png}
    \end{center}
    
    \begin{enumerate}
        \item[(1)] 图2中折线 ABC 表示\underline{\hspace{1.5cm}}槽中的深度与注水时间之间的关系，线段 DE 表示\underline{\hspace{1.5cm}}槽中的深度与注水时间之间的关系（以上两空均填“甲”或“乙”），点 B 的纵坐标表示的实际意义是\underline{\hspace{4cm}}。
        \vspace{1.5cm}
        \item[(2)] 注水多长时间时，甲、乙两个水槽中的水深相同？
        \vspace{2.5cm}
        \item[(3)] 若乙槽底面积为 36 平方厘米（壁厚不计），求乙槽中铁块的体积。
        \vspace{3.5cm}
    \end{enumerate}
\end{enumerate}

\vspace{0.8cm}

\noindent \textbf{四、勇敢闯一闯}

\begin{enumerate}
    \item[16.] (10分) 如图1，在长方形 ABCD 中，AB=12 厘米，BC=6 厘米。点 P 沿 AB 边从点 A 开始向点 B 以 2 厘米/秒的速度移动；点 Q 沿 DA 边从点 D 开始向点 A 以 1 厘米/秒的速度移动。如果 P、Q 同时出发，用 $t$ (秒) 表示移动的时间，那么：
    
    \begin{center}
        \includegraphics[width=0.45\textwidth]{16.png}
    \end{center}
    
    \begin{enumerate}
        \item[(1)] DQ=\underline{\hspace{2cm}}厘米，AP=\underline{\hspace{2cm}}厘米（用含 $t$ 的代数式表示）
        \vspace{1.5cm}
        \item[(2)] 如图 1，当 $t=$\underline{\hspace{2cm}}秒时，线段 AQ 与线段 AP 相等？
        \vspace{2.5cm}
        \item[(3)] 如图 2，P、Q 到达 B、A 后继续运动，P 点到达 C 点后都停止运动。当 $t$ 为何值时，线段 AQ 的长等于线段 CP 的长的一半。
        \vspace{3.5cm}
    \end{enumerate}


\end{enumerate}

\end{document}